\section{Survey Instrument and Measures}

\subsection{Survey Structure and Domains}

The survey will be structured into multiple domains covering the characteristics and experiences of respondents as members of the Fimfiction.net and wider My Little Pony community. Planned domains include demographic characteristics; disability and neurodivergence; fandom history and entry into the fandom; engagement with different generations of My Little Pony; engagement with Fimfiction.net; engagement with other online and offline fandom communities; Fimfiction group participation; fiction reading behaviour and preferences; character and pairing preferences; broader fandom interests; social relationships formed through the fandom; and creative participation, including fanfiction writing and use of generative artificial intelligence.

% TBD: The final ordering of domains, number of questions, branching or skip logic,
% and final survey length will be established during instrument development.

\subsection{Variables Collected}

The planned variables will cover demographic characteristics, fandom history and engagement, Fimfiction-specific activity, social and community participation, and creative participation.

Demographic variables will include age, country of residence, highest educational attainment, political orientation, religious affiliation, disability, neurodivergence, LGBTQ+ identity, and other relevant demographic characteristics. Fandom variables will include how and when respondents entered the My Little Pony fandom, age at fandom entry, generations engaged with, interest in G6, degree of fandom identification, and engagement with MLP content across other platforms. Fimfiction-specific variables will include site engagement, group participation, reading behaviour, story-rating preferences, tag preferences, character preferences, generation preferences, romance and pairing preferences, and engagement with original characters. Social and community variables will include convention attendance, online and offline friendships formed through the fandom, romantic relationships originating within the fandom, and participation in non-MLP activities with other fans. Writer-specific variables will include whether respondents write fanfiction and use of generative AI for activities such as grammar correction, proofreading, brainstorming, story development, prose generation, artwork, and other creative purposes.

% TBD: The final variable list and coding of each variable will be established during survey-instrument development.

\subsection{Sources or Adaptations of Survey Items}

\sourceimage{958760}{Ah know what Ah'm doin'!}

The survey will combine newly developed items specific to Fimfiction and My Little Pony fandom engagement with, where appropriate, items adapted from established instruments or existing research.

% TBD: The specific sources, instruments, individual adapted items, and procedures for evaluating
% their suitability, including any relevant permissions or licensing requirements, will be established
% before finalisation of the survey instrument.

\subsection{Response Scales}

The survey will use response formats appropriate to the characteristics of each variable. These will include categorical and multiple-choice responses, select-all-that-apply responses where multiple categories may apply, and numerical rating scales for preferences such as interest in Fimfiction tags and character categories. The proposed preference scale is a \textbf{1--10 scale}, where lower values represent low interest or active avoidance and higher values represent strong interest or active preference.

Questions concerning behaviours such as reading Mature-rated fiction, reading sexually explicit fiction, convention attendance, friendships, and platform engagement will use categorical frequency or participation responses as appropriate.

Select-all-that-apply responses will permit respondents to select multiple applicable categories, and each selected category will be treated independently for analysis. Percentages for such questions may therefore sum to more than 100\%.

% TBD: The precise response options, scale anchors, scale direction, and handling of
% ``prefer not to say,'' ``not applicable,'' and ``other'' responses will be established
% during finalisation of the survey instrument.

\subsection{Operational Definitions}

Fandom participation will be distinguished from broader identification with the fandom where appropriate, so that respondents can report both their behavioural engagement and whether they consider themselves a brony, pegasister, or another form of MLP fan.

Disability and neurodivergence will be measured as separate characteristics. Where respondents may legitimately have multiple applicable characteristics, the survey will permit multiple selections rather than requiring a single category.

% TBD: The final operational definitions and category lists for disability, neurodivergence, minority status,
% active Fimfiction engagement, fandom membership, writer status, and other important derived variables
% will be established before finalisation of the survey instrument.

\subsection{Construction of Composite and Derived Measures}

At present, the study does not require any predetermined composite measures. Most planned variables will be analysed individually.

Some variables may subsequently be combined or transformed for specific descriptive or exploratory analyses. Potential derived measures include the number of platforms used, number of fandoms engaged with, or number of disability or neurodivergence categories selected.

Raw responses will be retained wherever possible so that derived variables can be reproduced and alternative classifications can be evaluated.

% TBD: The final set of composite or derived measures, their source variables, calculation rules,
% handling of missing or non-applicable responses, and any thresholds used to create derived categories
% will be specified before the final analysis is conducted.
